% ======================================================================
% Closed-Loop Ant Colony Optimization
% ======================================================================

\documentclass[sigconf,anonymous,review]{acmart}

% , , , , , , , , , -
% Packages
% , , , , , , , , , -
\usepackage{amsmath, amssymb, amsfonts}
\usepackage{mathtools}
\usepackage{amsthm}
\usepackage{bm}

\theoremstyle{definition}
\newtheorem{definition}{Definition}[section]
\theoremstyle{plain}
\newtheorem{proposition}{Proposition}[section]
\newtheorem{lemma}[proposition]{Lemma}
\usepackage{booktabs}
\usepackage{graphicx}
\usepackage{subcaption}
\usepackage{multirow}
\usepackage{hyperref}
\usepackage{cleveref}
\usepackage{algorithm}
\usepackage{algorithmic}
\usepackage{xcolor}

% , , , , , , , , , -
% Convenience macros
% , , , , , , , , , -
\newcommand{\TODO}[1]{\textcolor{red}{[TODO: #1]}}
\newcommand{\eg}{e.g., }
\newcommand{\ie}{i.e., }
\newcommand{\etal}{et al.}
\newcommand{\R}{\mathbb{R}}
\newcommand{\cO}{\mathcal{O}}

% Problem names
\newcommand{\TSP}{\textsc{TSP}}
\newcommand{\CVRP}{\textsc{CVRP}}
\newcommand{\VRP}{\textsc{VRP}}
\newcommand{\ACO}{\textsc{ACO}}
\newcommand{\MMAS}{\textsc{MMAS}}
\newcommand{\LGA}{\textsc{LGA}}

% Box helper (no extra packages)
\newcommand{\DefBox}[1]{
\begin{center}
\fbox{\begin{minipage}{0.96\linewidth}
#1
\end{minipage}}
\end{center}
}

% , , , , , , , , , -
% Title / Author
% , , , , , , , , , -
\title{Closed-Loop Neural Control \\ of Ant Colony Optimization for Routing}

\author{Anonymous Author(s)}
\affiliation{
  \institution{Anonymous Institution}
}
\email{anonymous@domain.com}

\date{}

\begin{document}

% ======================================================================
% Abstract
% ======================================================================
\begin{abstract}
Ant Colony Optimization (ACO) often stagnates when pheromone values concentrate on suboptimal edges.
Existing neural-guided ACO methods provide \emph{static, instance-conditioned} guidance computed once before search, limiting their ability to react to the solver's evolving state.
We propose Learning-Guided ACO (\LGA{}), a \emph{closed-loop} controller that periodically observes pheromones and the incumbent solution, outputting per-edge residual logits that are held constant for a fixed number of inner iterations (sample-and-hold control).
The resulting transition kernel corresponds to exponential tilting of the baseline ACO policy and admits a KL-regularized interpretation.
Empirically, we find that \emph{localized local search} is crucial for stable policy-gradient training: it preserves sampled edges and improves action--outcome fidelity, whereas global local search frequently overwrites sampled decisions and destabilizes learning.
On Euclidean \TSP{} and \CVRP{} (via giant-tour + Split), \LGA{} improves over a strong perturbation-based ACO backend and open-loop guidance under matched wall-clock budgets.
\end{abstract}

\maketitle


% ======================================================================
% 1. Introduction
% ======================================================================
\section{Introduction}

Ant Colony Optimization (ACO) is an effective metaheuristic for routing problems, but it can stagnate: once pheromone values concentrate on a suboptimal solution, the algorithm loses the ability to explore alternatives.
Common remedies (restart, smoothing) are heuristic and do not condition on search state.
Recent neural ACO methods (e.g., DeepACO~\cite{deepaco}) address this by learning instance-conditioned guidance, but they compute this guidance \emph{once} before search begins.
Because such guidance is fixed throughout the run, it cannot respond when pheromones converge prematurely or improvements stall.

We propose \emph{Learning-Guided ACO} (\LGA{}), a closed-loop controller that periodically observes the solver state---pheromone distribution and incumbent solution---and outputs per-edge residual logits that are held constant for $S$ inner iterations (sample-and-hold).
This state-conditioned guidance enables the model to steer sampling toward edges that are promising \emph{given the current search trajectory}, not just the instance geometry.
The shaping corresponds to exponential tilting of the baseline ACO kernel and admits a KL-regularized interpretation (\Cref{sec:controller}).

We build on perturbation-based ACO~\cite{faco}, which bounds the number of edge changes per iteration, yielding a focused large-neighborhood search.
A key finding is that \emph{localized local search} (LS)---originally introduced for efficiency---is necessary for stable policy-gradient training: it preserves sampled edges (high action--outcome fidelity), whereas global LS overwrites them and weakens credit assignment (\Cref{sec:localized-ls}).

\DefBox{
\textbf{Open-loop vs.\ closed-loop neural ACO.}
Let $x$ be a problem instance and $\mathcal{S}_h$ the solver state at refresh step $h$ (pheromones and incumbent).
\emph{Open-loop}: guidance $g_\theta(\cdot; x)$ is computed once before search and remains fixed.
\emph{Closed-loop}: guidance $g_\theta(\cdot; \mathcal{S}_h)$ is recomputed at multiple refresh steps, conditioning on the evolving solver state.
}

\paragraph{Contributions.}
\begin{enumerate}
    \item \textbf{Closed-loop neural ACO.} We introduce a controller that conditions on pheromones and incumbent structure, refreshing guidance during search via sample-and-hold control (\Cref{sec:method}).
    \item \textbf{Key enabling insight.} We show that localized LS is necessary for stable policy-gradient training in our setting, and introduce edge survival as a diagnostic metric (\Cref{sec:localized-ls}).
    \item \textbf{CVRP extension.} We extend to CVRP via giant-tour representation with Split decoding (\Cref{sec:cvrp}).
    \item \textbf{Empirical validation.} Under matched wall-clock budgets, \LGA{} improves over both the vanilla backend and open-loop neural guidance on TSP and CVRP across scales (\Cref{sec:experiments}).
\end{enumerate}

% ======================================================================
% 2. Related Work
% ======================================================================
\section{Related Work}
\subsection{Neural Combinatorial Optimization}
Neural Combinatorial Optimization (NCO) applies deep learning to NP-hard routing problems~\cite{MLCO}.
\emph{End-to-end} methods generate solutions directly via attention-based decoding~\cite{pointernetwork,kool2019attention,kwon2020pomo}.
In contrast, \emph{Learning-Guided Optimization} (LGO) produces edge-level heuristics that guide a classical solver~\cite{RL4CO}, combining learned representations with principled search.

\subsection{Neural-Guided Ant Colony Optimization}
A prominent LGO line augments ACO with learned priors.
DeepACO~\cite{deepaco} trains a GNN to output edge logits that replace or augment the heuristic $\eta_{ij}$.
GFACS~\cite{gfacs} uses Generative Flow Networks; GTG-ACO~\cite{gtgaco} employs Graph Transformers.
Critically, all these methods compute guidance \emph{once} before search---they are \textbf{open-loop} and cannot adapt to evolving search dynamics.

\subsection{Perturbation-Based ACO}
A separate line improves ACO scalability by \textbf{reusing and partially modifying elite solutions}, making ACO behave like a large-neighborhood search (LNS) operator.
Focused ACO~\cite{faco} bounds the number of differences between a newly constructed solution and the incumbent, yielding focused search and tighter coupling with local search.
Related variants include PartialACO~\cite{partialaco}.
We treat perturbation-based ACO as a \textbf{scalable backend} and ask: \emph{can a neural controller improve this family further by reacting to solver state in closed loop?}

\subsection{Learning to Control Optimization}
Our work draws on \emph{learning to optimize}~\cite{bengio2021machine}, where neural networks parameterize components of optimization algorithms (e.g., branch-and-bound policies~\cite{gasse2019exact}, adaptive operator selection~\cite{li2019learning}).
A key distinction is that we implement \textbf{closed-loop feedback control}: the controller observes pheromone distribution and incumbent solution, then modulates future transitions accordingly.

\subsection{Positioning}
Table~\ref{tab:positioning} summarizes the landscape.
\LGA{} studies closed-loop control in which guidance conditions on solver state and is refreshed during the run.

\begin{table}[t]
\centering
\caption{Positioning of neural-guided ACO methods. \LGA{} builds on perturbation-based ACO~\cite{faco} as backend.}
\label{tab:positioning}
\resizebox{\columnwidth}{!}{%
\begin{tabular}{lcccc}
\toprule
Method & Guidance & Adapts to & Control & Scalable \\
    & Source   & Solver State? & Type & Backend \\
\midrule
DeepACO~\cite{deepaco} & GNN & \texttimes & Open-loop & Standard \\
GFACS~\cite{gfacs} & GFlowNet & \texttimes & Open-loop & Standard \\
GTG-ACO~\cite{gtgaco} & Transformer & \texttimes & Open-loop & Standard \\
\textbf{LGA (Ours)} & GNN + Context & \checkmark & Closed-loop & Perturbation \\
\bottomrule
\end{tabular}%
}
\end{table}

% ======================================================================
% 3. Background and Preliminaries
% ======================================================================
\section{Background and Preliminaries}
\label{sec:background}

\subsection{Problem Definitions}
We consider the Euclidean \TSP{} defined on a graph $G=(V, E)$ with $n = |V|$ nodes.
The goal is to find a Hamiltonian cycle $\pi$ minimizing the total tour length:
\begin{equation}
C(\pi) = \sum_{i=1}^{n-1} d(\pi_i, \pi_{i+1}) + d(\pi_n, \pi_1),
\end{equation}
where $d(\cdot, \cdot)$ denotes Euclidean distance.
For \CVRP{}, we additionally have a depot (node~0) and vehicle capacity constraint $Q$.
The objective is to find a set of routes, each starting and ending at the depot, that collectively visit all customers while respecting capacity limits.
We generate random instances with coordinates sampled uniformly from $[0, 1]^2$.

\paragraph{Candidate graph.}
To scale to large instances, we solve on a sparse $K$-nearest-neighbor candidate graph rather than the complete graph.
Tours are restricted to edges in this candidate set, which is standard practice in scalable \TSP{}/\CVRP{} heuristics~\cite{helsgaun2000effective}.
We use $K=32$ throughout.

\subsection{Ant Colony Optimization}
\label{sec:aco}
Ant Colony Optimization (ACO) constructs solutions probabilistically.
At each step, an ant at node $i$ chooses the next node $j$ from its feasible neighborhood $\mathcal{N}(i)$ using the transition rule:
\begin{equation}
p(j \mid i) = \frac{[\tau_{ij}]^\alpha [\eta_{ij}]^\beta}{\sum_{k \in \mathcal{N}(i)} [\tau_{ik}]^\alpha [\eta_{ik}]^\beta}, \quad j \in \mathcal{N}(i).
\label{eq:aco-transition}
\end{equation}
Here, $\tau_{ij}$ is the pheromone trail intensity, $\eta_{ij} = 1/d_{ij}$ is the heuristic visibility, and $\alpha, \beta$ control their relative influence.

\subsection{\MMAS{} and Stagnation}
\emph{Max-Min Ant System} (\MMAS{})~\cite{stutzle2000max} improves standard ACO by bounding pheromone values to $\tau_{ij} \in [\tau_{\min}, \tau_{\max}]$, preventing premature convergence to suboptimal solutions.
Only the iteration-best or global-best ant deposits pheromone, focusing reinforcement on high-quality solutions.
Despite these mechanisms, \emph{stagnation}---where $\tau$ converges to a suboptimal tour and exploration effectively ceases---remains a persistent challenge.

\subsection{Local Search}
Achieving competitive performance requires integrating local search (LS) into ACO pipelines.
For \TSP{}, 2-opt moves are standard.
While LS improves solution quality, it complicates learning by decoupling the policy's edge choices from the final reward---we address this through localized LS (\Cref{sec:localized-ls}).

% ======================================================================
% 4. Closed-Loop Ant Colony Optimization
% ======================================================================
% ======================================================================
% 4. Closed-Loop Controller
% ======================================================================
\section{Method}
\label{sec:method}

This section describes \LGA{}: a closed-loop neural controller for ACO.
We first give an overview (\Cref{sec:overview}), then detail the controller parameterization (\Cref{sec:state,sec:controller,sec:architecture}), training objective (\Cref{sec:training}), and pheromone stabilization (\Cref{sec:smoothed-mmas}).

\subsection{Overview}
\label{sec:overview}

We run a perturbation-based ACO backend~\cite{faco} on a sparse $K$-nearest-neighbor candidate graph.
Every $S$ inner iterations (one \emph{outer step} $h$):
\begin{enumerate}
    \item \textbf{Observe} the macro-state $\mathcal{S}_h = (x, \tau^{(h)}, b^{(h)})$: instance $x$, pheromones $\tau$, incumbent tour $b$.
    \item \textbf{Compute} shaping logits $z_\theta(\mathcal{S}_h) \in \R^{n \times K}$ via the neural controller.
    \item \textbf{Hold} these logits constant for $S$ inner iterations.
    \item \textbf{Sample} ants via bounded-difference construction (at most $M$ edge changes from incumbent), apply localized LS, update incumbent and pheromones.
\end{enumerate}
Total iterations: $T = H \cdot S$, where $H$ is the number of outer steps.
Pheromones are updated every inner iteration; we denote by $\tau^{(h)}$ the pheromone state at the \emph{start} of outer step $h$.

\paragraph{What the controller can and cannot do.}
The controller reweights candidate edges but cannot eliminate them (support preservation).
It operates within the candidate graph and bounded-perturbation constraints of the backend.
The backend handles pheromone updates, incumbent tracking, and localized LS.

\paragraph{Backend: perturbation-based construction.}
In each inner iteration, ants do not construct a tour from scratch.
Following~\cite{faco}, each ant starts from a copy of the incumbent tour and performs a constrained rewire procedure.
The procedure works as follows (see Algorithm~\ref{alg:perturbation} for pseudocode):

\begin{enumerate}
    \item \textbf{Initialize:} Copy the incumbent tour $b^{(h)}$ as the working tour. Select a random starting node $u_0$. Mark $u_0$ visited.
    \item \textbf{Sequential expansion:} Repeat until either (a) $M$ new edges have been created, or (b) a fixed step budget is exhausted:
    \begin{itemize}
        \item From current node $u$, compute transition probabilities over the candidate set $\mathcal{N}_K(u)$ (unvisited neighbors in the $K$-NN graph) using Eq.~\eqref{eq:held-logits}.
        \item Sample next node $v$ from this distribution.
        \item If edge $(u,v)$ is \emph{not} in the current working tour, it is a ``new edge.'' Increment the new-edge counter and add incident nodes to the LS checklist.
        \item \textbf{Relocate} $v$: remove $v$ from its current position and insert it as the successor of $u$. This maintains a valid tour at all times.
        \item Mark $v$ visited; set $u \leftarrow v$.
    \end{itemize}
    \item \textbf{Return:} The modified tour after at most $M$ new edges.
\end{enumerate}

The action space is thus a sequence of \emph{next-node choices} from the candidate graph, identical to classical ACO but constrained to terminate early.
The log-probability of a sampled tour is the sum of log-probabilities over stochastic node selections:
\begin{equation}
\log \pi_\theta(\pi_{a,s} \mid \mathcal{S}_h) = \sum_{t=1}^{T_a} \log p^\theta(v_t \mid u_t; h, s),
\label{eq:logp-trajectory}
\end{equation}
where $T_a \le M$ is the number of steps taken by ant $a$, and we sum only over stochastic decisions (steps with $>1$ feasible candidate).

\paragraph{Notation.}
We use $h = 1,\ldots,H$ for outer steps (controller refreshes) and $s = 1,\ldots,S$ for inner \ACO{} iterations between refreshes.
Total inner iterations: $T = H \cdot S$.
We write $b^{(h)}$ for the incumbent tour at outer step $h$, and $\tau^{(h,s)}$ for the pheromone state at inner iteration $s$ of outer step $h$.

\paragraph{Key design: sample-and-hold with evolving pheromones.}
The neural controller outputs shaping logits $z_\theta(\mathcal{S}_h)$ \emph{once} at the start of each outer step and holds them constant for $S$ inner iterations.
However, pheromones $\tau^{(h,s)}$ \emph{evolve} at every inner iteration via the standard \MMAS{} update rule.
At each inner iteration $s$, the transition kernel uses the \emph{current} pheromones $\tau^{(h,s)}$ combined with the \emph{held} neural logits $z_\theta(\mathcal{S}_h)$:
\begin{equation}
\log p^\theta(j \mid i; h, s) = \alpha \log \tau^{(h,s)}_{ij} + \beta \log \eta_{ij} + z_{ij}(\mathcal{S}_h) - \log Z_i.
\label{eq:held-logits}
\end{equation}
This design separates the \emph{fast} pheromone dynamics (updated every inner step) from the \emph{slow} neural control signal (updated every $S$ steps), analogous to sample-and-hold control in signal processing.

\subsection{Sample-and-Hold Control}
\label{sec:process}
The controller is invoked every $S$ inner steps.
At outer step $h$, it observes $\mathcal{S}_h = (x, \tau^{(h)}, b^{(h)})$ and outputs shaping logits $z_\theta(\mathcal{S}_h) \in \R^{n\times K}$, held constant for the next $S$ iterations.

\paragraph{Why sample-and-hold?}
Refreshing every iteration is expensive and potentially unstable; refreshing too rarely degenerates to open-loop.
Sample-and-hold trades off compute and adaptivity: the controller amortizes inference cost over $S$ sampling steps while still adapting as search progresses.
We ablate the $(H, S)$ factorization with fixed total budget in \Cref{sec:ablations}.


\subsection{State Representation}
\label{sec:state}
The controller conditions on the search context $\mathcal{S}_h$, which captures both static instance structure and dynamic search state.

\paragraph{Static features.}
A GNN encoder computes node embeddings $h_i \in \R^d$ from node coordinates and the $K$-NN distance graph.

\paragraph{Dynamic features.}
For each candidate edge $(i, j)$:
\begin{itemize}
    \item \emph{Relative pheromone:} $\log(\tau_{ij} / \bar{\tau}_i)$, where $\bar{\tau}_i$ is the mean pheromone over candidates of node $i$.
    \item \emph{Pheromone dispersion:} Coefficient of variation of pheromone values at node $i$.
    \item \emph{Incumbent structure:} Features indicating whether edge $(i,j)$ appears in the current source tour (see Appendix for details).
\end{itemize}

\subsection{Log-Space Shaping}
\label{sec:controller}
The controller outputs a scalar guidance score $z_{ij} = g_{\theta}(i, j; \mathcal{S}_h)$ for each candidate edge.
Let the baseline ACO logits be $\ell^{\text{ACO}}_{ij} = \alpha \log \tau^{(h)}_{ij} + \beta \log \eta_{ij}$.
The shaped transition kernel adds the learned logits in log-space:
\begin{equation}
\log p^\theta(j \mid i) = \ell^{\text{ACO}}_{ij} + z_{ij} - \log Z_i,
\label{eq:logspace-prior}
\end{equation}
where $Z_i$ is the partition function.

\paragraph{KL-regularized interpretation.}
Let $q(\cdot\mid i)$ be the baseline ACO distribution (without $z$).
Adding $z$ in log-space corresponds to exponential tilting:
$p^\theta(j\mid i) \propto q(j\mid i)\exp(z_{ij})$.
Equivalently, $p^\theta$ solves a KL-regularized objective:
\begin{equation}
p^\theta(\cdot\mid i) = \arg\max_{p \in \Delta(\mathcal{N}(i))}\Big(\mathbb{E}_{j\sim p}[z_{ij}] - \mathrm{KL}(p\,\|\,q)\Big).
\label{eq:kl-regularized}
\end{equation}
This makes the controller's role precise: it performs a KL-regularized policy improvement around the baseline kernel.

\paragraph{Support preservation.}
With $\tau_{\min} > 0$, $\eta_{ij} > 0$, and finite $z_{ij}$, we have $p^\theta(j \mid i) > 0$ for all candidates.
The controller reweights candidate edges but cannot eliminate them.

\subsection{Architecture}
\label{sec:architecture}
The controller comprises:
\begin{itemize}
    \item \textbf{Encoder}: A 12-layer GNN with residual connections and batch normalization, producing edge embeddings $w_{ij} \in \mathbb{R}^{32}$ from node coordinates and edge features.
    \item \textbf{Decoder}: A 3-layer MLP mapping edge embeddings to guidance logits $g_{\theta}(i, j)$.
\end{itemize}

\paragraph{What runs when.}
The GNN encoder is invoked once per outer step (every $S$ inner iterations), not every inner iteration.
The encoder receives both static features (node coordinates, candidate distances) and dynamic features (pheromone state, incumbent structure) as input.
Because these dynamic features change only at the outer-step boundary (we observe $\tau^{(h,0)}$ and $b^{(h)}$), the full forward pass amortizes over $S$ sampling steps.
The dominant cost is the GNN forward pass at $O(L \cdot n \cdot K \cdot d)$ for $L$ layers, $n$ nodes, $K$ candidates, and hidden dimension $d=32$.
This is invoked $H$ times per instance, yielding total neural cost $O(H \cdot L \cdot n \cdot K \cdot d)$ compared to $O(T \cdot m \cdot M \cdot K)$ for sampling (where $T = H \cdot S$, $m$ ants, and $M$ perturbations per ant).
With our default settings ($H=10$, $S=100$, $m=100$, $M=12$), the neural overhead is typically 5--15\% of total runtime at $n=10$K (see Table~\ref{tab:profiling}).

% ======================================================================
% 4.1 Smoothed MMAS
% ======================================================================
\subsection{Stabilized Pheromone Dynamics}
\label{sec:smoothed-mmas}
Standard \MMAS{} dynamically adjusts pheromone bounds based on the best tour cost, coupling the dynamics to objective scaling.
This coupling introduces instability when training neural controllers, as the input distribution (pheromone values) varies with instance difficulty.

\paragraph{What we change.}
We use \emph{fixed bounds} ($\tau_{\max}=1$, $\tau_{\min}=1/K$) independent of instance cost.
Let $B$ denote the edge set of the best solution (iteration-best or global-best).
We update each pheromone value via convex interpolation toward a fixed target:
\begin{equation}
\tau^{(h,s+1)}_{ij} \leftarrow (1-\rho)\tau^{(h,s)}_{ij} + \rho \cdot \mathcal{T}_{ij},
\label{eq:smoothed-mmas}
\end{equation}
where $\mathcal{T}_{ij} = \tau_{\max}$ if $(i,j) \in B$ and $\mathcal{T}_{ij} = \tau_{\min}$ otherwise.

\paragraph{Justification for $\tau_{\min} = 1/K$.}
In standard \MMAS{}, $\tau_{\min}$ is set so that even at maximum pheromone concentration, the probability of selecting a non-best edge remains above a threshold $p_{\text{best}}$.
With $K$ candidates per node and uniform $\tau_{\min}$, the minimum selection probability for any candidate is approximately $\tau_{\min} / (\tau_{\max} + (K-1)\tau_{\min})$.
Setting $\tau_{\min} = \tau_{\max}/K = 1/K$ ensures that even the least-favored candidate retains probability $\approx 1/(2K-1) \approx 1.6\%$ when $K=32$, preventing complete stagnation.
This choice is inspired by the classic \MMAS{} derivation but decouples the bounds from solution cost.

\paragraph{Why it helps learning.}
Because the bounds are fixed, the controller observes scale-stable inputs across instances of different difficulty.
The update rule~\eqref{eq:smoothed-mmas} is a contraction: pheromones remain in $[\tau_{\min}, \tau_{\max}]$ regardless of the best-set sequence, preventing the unbounded growth or collapse that can occur under cost-calibrated bounds.
This stabilizes training without sacrificing final solution quality (see ablation in Section~\ref{sec:ablations}).

\paragraph{Sensitivity to $\rho$ and $\tau_{\min}$.}
The evaporation rate $\rho$ controls how quickly pheromones adapt to new best solutions.
With fixed bounds, larger $\rho$ means faster convergence to the target but also faster forgetting of exploration history.
We use $\rho=0.1$ for TSP and $\rho=0.5$ for CVRP based on preliminary experiments.
The choice of $\tau_{\min}$ affects the exploration--exploitation tradeoff: smaller $\tau_{\min}$ allows stronger concentration but risks stagnation, while larger $\tau_{\min}$ maintains diversity at the cost of slower convergence.
Our default $\tau_{\min} = 1/K$ balances these concerns.
See Appendix~\ref{app:tuning} for sensitivity analysis.

% ======================================================================
% 4.2 Training Objective
% ======================================================================
\subsection{Learning Objective and Optimization}
\label{sec:training}
Let $\hat{\pi}_{a,s} = \mathrm{LS}(\pi_{a,s})$ denote the tour produced by ant $a$ at inner iteration $s$ after local search.
We optimize the expected per-iteration post-LS cost as a dense surrogate:
\begin{equation}
J(\theta) = \mathbb{E} \left[ \frac{1}{T} \sum_{h=1}^H \sum_{s=1}^S \frac{1}{m} \sum_{a=1}^m C(\hat{\pi}_{a,s}) \right].
\end{equation}

\paragraph{Why a dense objective?}
While evaluation typically reports best-so-far after a budget, we minimize average cost for three reasons:
(i) it provides lower-variance gradients than sparse best-so-far signals;
(ii) it empirically correlates with improved best-so-far under fixed budgets;
(iii) tracking incumbent quality separately confirms that better average translates to better final solutions.

We use the score function estimator (REINFORCE), differentiating through $\pi_{a,s}$ (pre-LS sampling):
\begin{multline}
\nabla_\theta J(\theta) \approx \frac{1}{T} \sum_{h,s} \frac{1}{m} \sum_{a=1}^m (C(\hat{\pi}_{a,s}) - b_s) \\
\times \nabla_\theta \log \pi_\theta(\pi_{a,s} \mid \mathcal{S}_h),
\end{multline}
where $b_s$ is the batch-mean baseline.
This estimator remains unbiased even with non-differentiable LS, as LS is part of the environment transition~\cite{williams1992simple}.

\paragraph{Trajectory structure.}
With perturbative sampling, each ant modifies at most $M$ edges.
Log-probabilities are summed over the realized edge choices.

\paragraph{Baseline and advantage.}
For each inner iteration $s$: $b_s = \frac{1}{m}\sum_{a=1}^{m} C(\hat{\pi}_{a,s})$.
Since we minimize cost, we define advantage as $A_{a,s} = b_s - C(\hat{\pi}_{a,s})$ (positive when better than average).
For PPO, we normalize advantages: $\hat{A}_{a,s} = (A_{a,s} - \mu_A) / (\sigma_A + \epsilon)$.

\paragraph{PPO-based Optimization.}
We implement a PPO variant for improved sample efficiency.
After collecting trajectories for $S$ inner iterations under $\pi_{\theta_{\text{old}}}$, we perform $K$ optimization epochs using the clipped surrogate:
\begin{multline}
\mathcal{L}^{\text{PPO}}(\theta) = -\mathbb{E}_{a,s}\Big[ \min\big( r_{a,s}(\theta) \hat{A}_{a,s}, \\
\text{clip}(r_{a,s}(\theta), 1{\pm}\epsilon) \hat{A}_{a,s} \big) \Big],
\end{multline}
where $r_{a,s}(\theta) = \pi_\theta(\pi_{a,s} | \mathcal{S}_h) / \pi_{\theta_{\text{old}}}(\pi_{a,s} | \mathcal{S}_h)$, $\epsilon=0.2$, and $\hat{A}_{a,s}$ is the normalized advantage (positive for better-than-average solutions).

\paragraph{What gets stored for PPO replay.}
To compute the probability ratio $r_{a,s}(\theta)$, we must replay the log-probability of each sampled trajectory under both old and new parameters.
Since the transition kernel~\eqref{eq:held-logits} depends on $\tau^{(h,s)}$ (which evolves during rollout), we store for each inner iteration $s$:
(i) the pheromone snapshot $\tau^{(h,s)}$ at sampling time,
(ii) the trajectory trace (sequence of node choices and which were stochastic), and
(iii) the post-LS cost $C(\hat{\pi}_{a,s})$.
During PPO epochs, we recompute $z_\theta(\mathcal{S}_h)$ with current parameters and replay log-probabilities using the stored pheromone snapshots, ensuring consistent probability ratio estimation.

\paragraph{Prior annealing.}
Optionally, we anneal the neural prior's influence within each control period, decaying the scale factor from $\gamma$ to $\gamma_{\min}$ across inner iterations.
This encourages stronger neural guidance early, gradually deferring to pheromone-driven dynamics.

\begin{algorithm}[t]
\caption{PPO Training for \LGA{}}
\label{alg:train}
\begin{algorithmic}[1]
\STATE Initialize policy $\theta$, optimizer
\FOR{epoch $e = 1, \dots, E$}
    \FOR{step $= 1, \dots, \text{Steps}$}
        \STATE Sample random instance $x$
        \STATE Initialize ACO state
        \FOR{outer step $h = 1$ to $H$}
            \STATE Build features $\mathcal{S}_h$; compute $z_{\theta_{\text{old}}}(\mathcal{S}_h)$
            \STATE \textbf{// Rollout: collect trajectories}
            \FOR{inner step $s = 1$ to $S$}
                \STATE Sample $m$ ants $\pi_{a,s}$; store $\log \pi_{\theta_{\text{old}}}(\pi_{a,s} | \mathcal{S}_h)$, $\tau^{(h)}$
                \STATE Apply localized LS: $\hat{\pi}_{a,s} \leftarrow \mathrm{LS}(\pi_{a,s})$; compute $C(\hat{\pi}_{a,s})$
                \STATE Update pheromone and incumbent
            \ENDFOR
            \STATE \textbf{// PPO update}
            \FOR{PPO epoch $= 1$ to $K$}
                \STATE Recompute $z_\theta(\mathcal{S}_h)$ with current parameters
                \FOR{inner step $s = 1$ to $S$}
                    \STATE Compute $\log \pi_\theta(\pi_{a,s} | \mathcal{S}_h)$ using stored $\tau^{(h)}$
                    \STATE Ratio: $r_{a,s} \leftarrow \exp(\log \pi_\theta - \log \pi_{\theta_{\text{old}}})$
                    \STATE Advantage: $\hat{A}_{a,s} \leftarrow \text{normalize}(\bar{C}_s - C(\hat{\pi}_{a,s}))$ \COMMENT{positive = better}
                    \STATE Clipped loss: $\mathcal{L}_s \leftarrow -\min(r_{a,s} \hat{A}_{a,s}, \text{clip}(r_{a,s}, 1{\pm}\epsilon) \hat{A}_{a,s})$
                \ENDFOR
                \STATE Backprop $\sum_s \mathcal{L}_s$; optimizer step (clip grad norm 1.0)
            \ENDFOR
        \ENDFOR
    \ENDFOR
\ENDFOR
\end{algorithmic}
\end{algorithm}

\begin{algorithm}[t]
\caption{Perturbation-Based Ant Construction}
\label{alg:perturbation}
\begin{algorithmic}[1]
\STATE \textbf{Input:} Incumbent tour $b$, starting node $u_0$, max new edges $M$, shaped kernel $p^\theta$
\STATE $\text{tour} \leftarrow \text{copy}(b)$; $\text{visited}[u_0] \leftarrow \text{true}$; $\text{new\_edges} \leftarrow 0$
\STATE $\text{checklist} \leftarrow \{u_0\}$; $u \leftarrow u_0$; $\log p \leftarrow 0$
\WHILE{$\text{new\_edges} < M$ \textbf{and} not all nodes visited}
    \STATE $\mathcal{C} \leftarrow \{v \in \mathcal{N}_K(u) : \neg\text{visited}[v]\}$ \COMMENT{feasible candidates}
    \IF{$|\mathcal{C}| = 0$}
        \STATE $v \leftarrow$ fallback to backup list or greedy
    \ELSIF{$|\mathcal{C}| = 1$}
        \STATE $v \leftarrow$ the single candidate \COMMENT{deterministic, $\log p$ unchanged}
    \ELSE
        \STATE Sample $v \sim p^\theta(\cdot \mid u)$ restricted to $\mathcal{C}$
        \STATE $\log p \leftarrow \log p + \log p^\theta(v \mid u)$
    \ENDIF
    \IF{$(u, v) \notin \text{edges}(\text{tour})$}
        \STATE $\text{new\_edges} \leftarrow \text{new\_edges} + 1$
        \STATE $\text{checklist} \leftarrow \text{checklist} \cup \{u, v\}$
    \ENDIF
    \STATE \textsc{Relocate}(tour, $u$, $v$) \COMMENT{move $v$ to be successor of $u$}
    \STATE $\text{visited}[v] \leftarrow \text{true}$; $u \leftarrow v$
\ENDWHILE
\STATE \textbf{Return:} tour, checklist, $\log p$
\end{algorithmic}
\end{algorithm}

% ======================================================================
% 5. Extended Analysis: Perturbative Search
% ======================================================================
\section{Key Insight: Localized LS for Stable Learning}
\label{sec:perturbative}

We adopt the local search procedure from~\cite{faco}, which restricts search moves to neighborhoods induced by the modified edges (rather than performing global 2-opt over the entire tour).
While this restriction was originally introduced for scalability, we find it is \emph{critical for stable policy-gradient training in our implementation}.

\subsection{Localized Local Search}
\label{sec:localized-ls}

\DefBox{
\textbf{Key Observation.} In our setting, global LS weakens credit assignment by overwriting a large fraction of sampled edges; localized LS preserves sampled decisions and stabilizes policy gradients. We quantify this via \emph{edge survival} and gradient variance (\Cref{sec:ablations}).
}

The procedure works as follows:
\begin{enumerate}
    \item \emph{Tracking:} We maintain a \texttt{checklist} of nodes incident to the $M$ new edges.
    \item \emph{Focused Optimization:} We apply restricted 2-opt only to edges involving nodes in the \texttt{checklist}.
\end{enumerate}
This yields $O(M \cdot K)$ move evaluations vs.\ $O(N \cdot K)$ for global LS.

\paragraph{Credit assignment fidelity.}
Let $\pi$ denote the pre-LS tour sampled by an ant, and let $\hat{\pi}=\mathrm{LS}(\pi)$ be the post-LS tour that determines the training signal.
Under global LS, $\hat{\pi}$ can differ from $\pi$ in $O(N)$ edges even when the ant only made $O(M)$ targeted modifications.
The final cost $C(\hat{\pi})$ may then depend weakly on the specific edges sampled under the controller, leading to high gradient variance.
In contrast, localized LS confines modifications to a neighborhood around the ant's changed edges, increasing the overlap between sampled actions and the resulting solution.
We measure this via \emph{edge survival}: $|E(\pi)\cap E(\hat{\pi})|/|E(\pi)|$.

\paragraph{Budget matching.}
To ensure fair comparison between localized and global LS, our ablation (Table~\ref{tab:ablation-ls}) uses the same number of 2-opt iterations for both variants.
Specifically, global LS is given the same number of node-scanning passes as localized LS (not more LS power, just different scope).
The key difference is \emph{which} edges are examined: localized LS focuses on the $O(M \cdot K)$ edges incident to the checklist, while global LS scans $O(N \cdot K)$ edges but with the same iteration budget.

% ======================================================================
% 6. Extension to CVRP
% ======================================================================
\section{Extension to \CVRP{}}
\label{sec:cvrp}
To demonstrate the generality of our framework, we extend \LGA{} to the Capacitated Vehicle Routing Problem.

\paragraph{Giant-tour representation.}
We represent a CVRP solution by a \textbf{giant tour} (a permutation of customers) and recover depot insertions via a Split dynamic program~\cite{prins2004}.
This keeps the search space in the same \emph{permutation domain} as TSP---so we can reuse the same perturbative construction, candidate graph, and localized 2-opt---while evaluating the true CVRP objective through Split after each ant.

\paragraph{Proxy mismatch.}
Improving Euclidean giant-tour length is not guaranteed to improve post-split CVRP cost.
However, Split depends heavily on customer order; improving geometric order typically reduces inter-customer travel and often improves Split cost.
The proxy is used only within LS---the reward signal for policy optimization is always the true CVRP cost computed via Split.

\subsection{Giant Tour Representation}
The \emph{Giant Tour} representation~\cite{prins2004} encodes a \CVRP{} solution as a permutation of all $n$ customers.
Given a permutation, the optimal route partition (depot insertions) is computed via the \texttt{Split} algorithm---a dynamic programming procedure that finds the minimum-cost way to segment the sequence into capacity-feasible routes.

For efficiency, we implement Split in C++ with prefix sums, enabling $O(1)$ cost evaluation for 2-opt moves.
Full algorithmic details appear in Appendix~\ref{app:split}.

\subsection{Modifications}
\begin{itemize}
    \item \emph{Graph:} The graph contains $n$ customer nodes. The depot is implicit in the objective evaluations.
    \item \emph{Feasibility Masking:} Identical to TSP (Hamiltonian path on customers).
    \item \emph{Localized LS:} The 2-opt operator optimizes the Euclidean length of the giant tour as a proxy. The true CVRP cost is computed only once per ant after the LS phase.
\end{itemize}

% ======================================================================
% 7. Experiments
% ======================================================================
\section{Experiments}
\label{sec:experiments}

We evaluate \LGA{} on random Euclidean \TSP{} and \CVRP{} instances across a range of problem sizes.
Our experiments address three questions:
(1) Does closed-loop control improve upon open-loop neural guidance?
(2) What is the contribution of each component (dynamic features, Smoothed-MMAS, localized LS)?
(3) How does \LGA{} scale to large instances?

\subsection{Experimental Setup}
We evaluate on instances of size $N \in \{1\text{K}, 5\text{K}, 10\text{K}, 50\text{K}, 100\text{K}\}$.
Training data is generated on-the-fly (uniform random distribution); validation uses fixed datasets of 128 instances.
All experiments use a single GPU (NVIDIA RTX 4090).

\paragraph{Two evaluation regimes.}
We define two regimes to ensure fair and realistic comparisons:
\begin{itemize}
    \item \emph{Standard regime} ($N \le 1\text{K}$): Compare against DeepACO, GFACS, GTG-ACO, and classical solvers. LKH-3/HGS serve as optimality references.
    \item \emph{Large-scale regime} ($N \ge 5\text{K}$): Compare against scalable methods only (vanilla MFACO, open-loop neural, FocusedACO). We report \emph{relative improvement} vs.\ the strongest scalable baseline under matched time, since LKH-3/HGS may not run reliably or within reasonable time at these scales.
\end{itemize}

\emph{ACO Hyperparameters:} $\alpha=1.0$, $\beta=1.0$, $\rho=0.1$ (\TSP{}) / $\rho=0.5$ (\CVRP{}).
Ants $m=100$. Candidate list $K=32$.
Perturbation limit $M=12$ (\TSP{}) / $M=16$ (\CVRP{}).
Outer iterations $H=10$, inner iterations $S=100$.

\emph{Training Hyperparameters:} PPO with $\epsilon=0.2$, $K=4$ PPO epochs per outer iteration.
Learning rate $5 \times 10^{-6}$ with cosine annealing.
Prior annealing from $\gamma=1.0$ to $\gamma_{\min}=0.2$.

\subsection{Metrics}
In addition to optimality gap, we track metrics to quantify the neural controller's behavior:
\begin{itemize}
    \item \emph{Gap (\%):} Relative gap to LKH-3 (TSP) or HGS (CVRP).
    \item \emph{Edge survival:} Fraction of sampled edges remaining after LS: $|E(\pi)\cap E(\hat{\pi})|/|E(\pi)|$. Higher survival means better action--outcome fidelity.
    \item \emph{KL divergence:} $\mathrm{KL}(p^\theta_h \| q_h)$ between shaped and baseline distributions, measuring control effort.
    \item \emph{Turnover:} Fraction of top-5\% probability edges that change between outer steps.
\end{itemize}

\subsection{Baselines}
We compare against:
\begin{enumerate}
    \item \emph{Classical Solvers:}
    \begin{itemize}
        \item \emph{Vanilla ACO:} Our perturbation-based ACO backend without the neural controller.
        \item \emph{LKH-3 / HGS:} Reference solvers (standard regime only; used for gap computation).
    \end{itemize}
    \item \emph{Neural-\ACO{} Baselines (standard regime only):}
    \begin{itemize}
        \item \emph{DeepACO}~\cite{deepaco}: Open-loop neural prior.
        \item \emph{GFACS}~\cite{gfacs}: Generative Flow Networks.
        \item \emph{GTG-ACO}~\cite{gtgaco}: Graph Transformer Guided.
    \end{itemize}
    \item \emph{Open-Loop Baseline:} Our model computed once from static features, to isolate the benefit of closed-loop feedback (all scales).
\end{enumerate}

% ,  NEW: fairness / budget subsection (reviewer-critical) , 
\subsection{Budget Matching and Fair Evaluation Protocol}
\label{sec:fairness}
To ensure fair comparison, we report results under \emph{matched wall-clock time}.

\paragraph{Runtime measurement.}
All experiments run on a single NVIDIA RTX 4090 GPU with 16 CPU cores (AMD Ryzen 9).
We measure wall-clock time from instance loading to final solution, including all overhead (feature extraction, GNN inference, sampling, LS, pheromone updates).
GPU time is included when applicable (GNN forward/backward passes).
We exclude JIT compilation warm-up by running a dummy instance before timing.

\paragraph{Implementation parity.}
All methods (LGA, vanilla ACO, open-loop neural) use the same C++ backend for sampling and local search, ensuring that performance differences reflect algorithmic choices rather than implementation quality.
The only difference is whether neural guidance is applied (and if so, how often it is refreshed).

\paragraph{Candidate graph restriction.}
All methods operate on the same $K$-nearest-neighbor candidate graph.
This is standard practice for scalable TSP/CVRP solvers and ensures fair comparison at large scales.

\paragraph{Baseline iteration adjustment.}
For vanilla ACO, we adjust the number of iterations so that total wall-clock time matches LGA (which includes GNN overhead).
For example, if LGA runs $H=10$ outer steps with $S=100$ inner iterations (1000 total) but takes 10\% longer due to neural inference, we give vanilla ACO approximately 1100 iterations to match runtime.

\subsection{Main Results}
\label{sec:main-results}

\paragraph{Experiment A1: TSP Results (Table~\ref{tab:main-tsp}).}
We evaluate \TSP{} performance across all scales.
For standard regime ($N=1\text{K}$), we report gap (\%) relative to LKH-3.
For large-scale regime ($N \ge 5\text{K}$), we report gap vs.\ best-known where available, and relative improvement (\%) vs.\ vanilla MFACO under matched wall-clock time.

\begin{table}[t]
\centering
\caption{\TSP{} Results. Standard regime: gap (\%) to LKH-3. Large-scale regime: gap to best-known and relative improvement vs.\ vanilla ACO under matched time (60s/instance for $N=5$K, 300s for $N=10$K, 1800s for $N\ge 50$K).}
\label{tab:main-tsp}
\small
\begin{tabular}{lccccc}
\toprule
Method & 1K & 5K & 10K & 50K & 100K \\
\midrule
\multicolumn{6}{l}{\emph{Standard Regime (Gap \% to LKH-3)}} \\
Vanilla ACO & -- & & & & \\
Open-Loop Neural & -- & & & & \\
\textbf{LGA (Ours)} & -- & & & & \\
\midrule
\multicolumn{6}{l}{\emph{Large-Scale Regime (Rel.\ Improv.\ \% vs.\ Vanilla)}} \\
Vanilla ACO & & 0.0 & 0.0 & 0.0 & 0.0 \\
Open-Loop Neural & & -- & -- & -- & -- \\
\textbf{LGA (Ours)} & & -- & -- & -- & -- \\
\bottomrule
\end{tabular}
\end{table}

\paragraph{Experiment A2: CVRP Results (Table~\ref{tab:main-cvrp}).}
Analogous evaluation for \CVRP{} with HGS as reference for standard regime.
We use capacity $Q=50$ for all scales.

\begin{table}[t]
\centering
\caption{\CVRP{} Results. Standard regime: gap (\%) to HGS. Large-scale regime: relative improvement vs.\ vanilla baseline. Time budgets matched per scale.}
\label{tab:main-cvrp}
\small
\begin{tabular}{lccccc}
\toprule
Method & 1K & 5K & 10K & 50K & 100K \\
\midrule
\multicolumn{6}{l}{\emph{Standard Regime}} \\
Vanilla ACO & -- & & & & \\
\textbf{LGA (Ours)} & -- & & & & \\
\midrule
\multicolumn{6}{l}{\emph{Large-Scale Regime (Rel.\ Improv.\ \%)}} \\
Vanilla ACO & & 0.0 & 0.0 & 0.0 & 0.0 \\
\textbf{LGA (Ours)} & & -- & -- & -- & -- \\
\bottomrule
\end{tabular}
\end{table}

\paragraph{Experiment A3: Comparison with Neural-\ACO{} Baselines (Table~\ref{tab:neural-baselines}).}
At $N=1\text{K}$ (standard regime, where prior neural-\ACO{} methods are computationally feasible), we compare against DeepACO, GFACS, and GTG-ACO under matched time budgets (30s/instance).

\begin{table}[t]
\centering
\caption{Comparison with neural-\ACO{} baselines at $N=1\text{K}$ (standard regime).}
\label{tab:neural-baselines}
\resizebox{\columnwidth}{!}{%
\begin{tabular}{lcccc}
\toprule
Method & \TSP{} Gap(\%) & \TSP{} Time(s) & \CVRP{} Gap(\%) & \CVRP{} Time(s) \\
\midrule
DeepACO & -- & -- & -- & -- \\
GFACS & -- & -- & -- & -- \\
GTG-ACO & -- & -- & -- & -- \\
\textbf{LGA (Ours)} & -- & -- & -- & -- \\
\bottomrule
\end{tabular}%
}
\end{table}

\paragraph{Experiment A4: Convergence Dynamics (Figure~\ref{fig:convergence}).}
We plot gap vs.\ wall-clock time curves showing how LGA converges faster than baselines.

\begin{figure}[t]
    \centering
    \begin{subfigure}{0.48\textwidth}
        \includegraphics[width=\linewidth]{gap_vs_h.pdf}
        \caption{Gap vs.\ Time: \TSP{}-10K.}
        \label{fig:gap-tsp}
    \end{subfigure}
    \begin{subfigure}{0.48\textwidth}
        \includegraphics[width=\linewidth]{prior_change_metrics.pdf}
        \caption{Gap vs.\ Time: \CVRP{}-10K.}
        \label{fig:gap-cvrp}
    \end{subfigure}
    \caption{Convergence curves: LGA (solid), vanilla ACO (dashed), open-loop neural (dotted). Shaded regions: $\pm 1$ std over 5 seeds.}
    \label{fig:convergence}
\end{figure}

\paragraph{Experiment A5: Active Control Verification (Figure~\ref{fig:active-control}).}
To verify that the controller is actively adapting rather than providing static guidance, we track two diagnostic metrics across outer steps:
\begin{itemize}
    \item \emph{KL divergence:} $\mathrm{KL}(p^\theta_h \| q_h)$ between the shaped distribution and the baseline ACO distribution (without neural prior) at each refresh $h$.
    \item \emph{Turnover:} Fraction of top-5\% probability edges that change between consecutive outer steps.
\end{itemize}
Figure~\ref{fig:active-control} shows that both metrics remain non-zero throughout the run, confirming that the controller continuously adjusts guidance in response to evolving pheromones and incumbent structure---not merely providing a static prior.

\begin{figure}[t]
    \centering
    \includegraphics[width=0.9\columnwidth]{active_control_metrics.pdf}
    \caption{Active control diagnostics on TSP-10K. Left: KL divergence between shaped and baseline distributions across outer steps. Right: Turnover (fraction of top-5\% edges changed between refreshes). Non-zero values throughout indicate the controller is actively adapting, not providing static guidance.}
    \label{fig:active-control}
\end{figure}

% ======================================================================
% Ablations
% ======================================================================
\subsection{Ablation Studies}
\label{sec:ablations}
We validate the core design choices of \LGA{} through four focused ablations, each addressing a specific research question.
All experiments use $N=1\text{K}$ with 5 seeds; hyperparameter and RL algorithm sweeps are deferred to Appendix~\ref{app:tuning}.

% ----------------------------------------------------------------------
% Ablation 1: Open-loop vs Closed-loop (Refresh Frequency)
% ----------------------------------------------------------------------
\paragraph{Ablation 1: Open-loop vs.\ Closed-loop Control (Table~\ref{tab:ablation-refresh}).}
\emph{Question: Does refreshing guidance during search improve over static, open-loop guidance?}

We vary the number of outer steps $H$ while fixing total iterations $T = H \times S = 1000$.
$H=1$ corresponds to open-loop control (guidance computed once); larger $H$ enables more frequent adaptation.

\begin{table}[t]
\centering
\caption{Open-loop vs.\ closed-loop control. Increasing refresh frequency ($H$) improves solution quality up to a point; excessive refreshes add overhead without benefit.}
\label{tab:ablation-refresh}
\small
\begin{tabular}{ccccc}
\toprule
$H$ & $S$ & Control Type & \TSP{} Gap (\%) & Overhead (\%) \\
\midrule
1 & 1000 & Open-loop & -- & -- \\
5 & 200 & Closed-loop & -- & -- \\
10 & 100 & Closed-loop & -- & -- \\
20 & 50 & Closed-loop & -- & -- \\
50 & 20 & Closed-loop & -- & -- \\
\bottomrule
\end{tabular}
\end{table}

\emph{Finding:} Closed-loop control ($H \ge 5$) consistently outperforms open-loop ($H=1$), confirming that adapting to solver state provides value beyond static instance-conditioned guidance.

% ----------------------------------------------------------------------
% Ablation 2: Dynamic Features
% ----------------------------------------------------------------------
\paragraph{Ablation 2: Dynamic Feature Removal (Table~\ref{tab:ablation-features}).}
\emph{Question: Which dynamic features drive the closed-loop benefit?}

We progressively remove dynamic features from the state representation to isolate their contributions.

\begin{table}[t]
\centering
\caption{Impact of dynamic features. Removing search context degrades performance toward open-loop levels.}
\label{tab:ablation-features}
\small
\begin{tabular}{lcc}
\toprule
Features & \TSP{} Gap (\%) & $\Delta$ vs.\ Full \\
\midrule
\textbf{Full (all dynamic)} & -- & -- \\
$-$ Incumbent structure & -- & -- \\
$-$ Pheromone dispersion & -- & -- \\
$-$ Relative pheromone & -- & -- \\
Static only (coords) & -- & -- \\
\bottomrule
\end{tabular}
\end{table}

\emph{Finding:} Pheromone features contribute most; incumbent structure provides additional signal. Static-only degenerates to open-loop performance.

% ----------------------------------------------------------------------
% Ablation 3: Localized vs Global LS
% ----------------------------------------------------------------------
\paragraph{Ablation 3: Localized vs.\ Global Local Search (Table~\ref{tab:ablation-ls}).}
\emph{Question: Why is localized LS necessary for stable learning?}

We compare localized LS (FACO-style, restricted to perturbed edges) against global 2-opt (unrestricted), measuring solution quality, edge survival (fraction of sampled edges retained after LS), and gradient variance during training.

\begin{table}[t]
\centering
\caption{Localized vs.\ global LS. Global LS overwrites sampled edges, reducing action--outcome fidelity and destabilizing learning.}
\label{tab:ablation-ls}
\small
\begin{tabular}{lccc}
\toprule
LS Type & Gap (\%) & Survival & Grad Var \\
\midrule
Global 2-opt & -- & -- & -- \\
Localized (FACO) & -- & -- & -- \\
\bottomrule
\end{tabular}
\end{table}

\emph{Finding:} Global LS achieves low survival ($<30\%$), causing high gradient variance and unstable training. Localized LS preserves $>80\%$ of sampled edges, enabling stable policy improvement.

% ----------------------------------------------------------------------
% Ablation 4: Smoothed MMAS
% ----------------------------------------------------------------------
\paragraph{Ablation 4: Stabilized vs.\ Standard Pheromone Bounds (Table~\ref{tab:ablation-mmas}).}
\emph{Question: Does stabilizing pheromone dynamics aid learning?}

Standard \MMAS{} adjusts $\tau_{\min}, \tau_{\max}$ based on best-tour cost, coupling pheromone scale to instance difficulty. Our stabilized variant uses fixed bounds ($\tau_{\max}=1$, $\tau_{\min}=1/K$).

\begin{table}[t]
\centering
\caption{Pheromone bound strategies. Fixed bounds stabilize training without sacrificing solution quality.}
\label{tab:ablation-mmas}
\small
\begin{tabular}{lccc}
\toprule
Bound Type & Gap (\%) & Train Loss Var & Pheromone CV \\
\midrule
Standard \MMAS{} & -- & -- & -- \\
Stabilized (fixed) & -- & -- & -- \\
\bottomrule
\end{tabular}
\end{table}

\emph{Finding:} Fixed bounds reduce training loss variance while achieving comparable or better final gaps, validating the stabilization design.

% ======================================================================
% Runtime and Scalability
% ======================================================================
\subsection{Runtime and Scalability}
\label{sec:runtime}

\paragraph{Experiment D1: Profiling Breakdown (Table~\ref{tab:profiling}).}
We decompose wall-clock time into components: GNN inference, ant sampling, local search, and pheromone updates.

\begin{table}[t]
\centering
\caption{Runtime profiling breakdown per outer iteration at different scales.}
\label{tab:profiling}
\small
\begin{tabular}{lcccc}
\toprule
$N$ & GNN (\%) & Sampling (\%) & LS (\%) & Update (\%) \\
\midrule
1K & -- & -- & -- & -- \\
10K & -- & -- & -- & -- \\
100K & -- & -- & -- & -- \\
\bottomrule
\end{tabular}
\end{table}

\paragraph{Experiment D2: Scaling Curve (Figure~\ref{fig:scaling}).}
We plot gap vs.\ problem size under fixed time budgets, comparing LGA, vanilla ACO, and prior work.

\begin{figure}[t]
    \centering
    \begin{subfigure}{0.48\textwidth}
        \includegraphics[width=\linewidth]{scaling_tsp.pdf}
        \caption{\TSP{} scaling.}
        \label{fig:scaling-tsp}
    \end{subfigure}
    \begin{subfigure}{0.48\textwidth}
        \includegraphics[width=\linewidth]{scaling_cvrp.pdf}
        \caption{\CVRP{} scaling.}
        \label{fig:scaling-cvrp}
    \end{subfigure}
    \caption{Scaling curves: Gap (\%) vs.\ problem size $N$ under matched time budgets. LGA maintains competitive gaps at large scales where prior neural-ACO methods become infeasible.}
    \label{fig:scaling}
\end{figure}

\paragraph{Experiment D3: Neural Overhead Analysis.}
The neural navigator adds overhead per outer iteration. We report the percentage overhead relative to vanilla MFACO and show that faster convergence compensates, yielding better solutions within the same total wall-clock time.

% ======================================================================
% 8. Discussion, Limitations, and Future Work
% ======================================================================
\section{Discussion and Limitations}
\LGA{} demonstrates that conditioning on solver state---pheromone distributions and incumbent structure---provides actionable information that static instance features cannot capture.
Our ablations show that localized LS improves edge survival over global LS and reduces gradient variance, enabling stable policy-gradient training.

\paragraph{Limitations.}
Our approach incurs overhead from the GNN encoder.
The RL training requires careful hyperparameter tuning, particularly for the control period $S$ and prior annealing schedule.
The perturbation-based sampling scheme, while enabling scalability, restricts reachability compared to global construction---this is an inherited trade-off from the backend design~\cite{faco}.

\paragraph{Future Work.}
Promising directions include: (i) applying closed-loop neural control to other combinatorial domains such as Job Shop Scheduling; (ii) exploring Transformer architectures for richer context encoding; and (iii) investigating curriculum learning strategies that gradually increase problem difficulty during training.

% ,  NEW: reproducibility checklist section (often appreciated) , 
\section{Reproducibility Checklist}
\label{sec:repro}
\begin{itemize}
    \item \emph{Code:} Training and evaluation scripts are included in this repository; we plan to release a cleaned version upon publication.
    \item \emph{Hyperparameters:} Candidate list $K=32$, ants $m=100$, $H=10$ outer iterations, $S=100$ inner iterations (\texttt{mini\_H}). \TSP{} uses $\rho=0.1$, $M=12$; \CVRP{} uses $\rho=0.5$, $M=16$. Learning rate defaults are PPO: $5\times 10^{-6}$, REINFORCE: $10^{-4}$ (see unified training script). Optimizer: AdamW with cosine annealing.
    \item \emph{Model:} 12-layer GNN encoder (hidden dim 32) and a 3-layer MLP edge decoder implemented in PyTorch/PyG.
\end{itemize}

% ======================================================================
% 9. Conclusion
% ======================================================================
\section{Conclusion}
We introduced \LGA{}, a framework for closed-loop neural control of ACO.
By conditioning guidance on solver state (pheromones and incumbent) and refreshing it during search, the controller can steer sampling in response to optimization dynamics.
Combined with a scalable perturbation-based backend and stabilized pheromone updates, \LGA{} outperforms both vanilla ACO and open-loop neural guidance on \TSP{} and \CVRP{} under matched budgets.
A key finding is that localized LS, originally designed for efficiency, is necessary for stable policy-gradient training due to improved action--outcome fidelity.
This work opens a path toward more adaptive optimization systems that learn to navigate complex search landscapes.

% ======================================================================
% Acknowledgments
% ======================================================================
\section*{Acknowledgments}
(Omitted for review)

% ======================================================================
% References
% ======================================================================
\bibliographystyle{plain}
\bibliography{references}

% ======================================================================
% Appendix
% ======================================================================
\appendix

\section{Additional Implementation Details}
\emph{Model Architecture:}
\begin{itemize}
    \item \emph{Encoder:} 12-layer GNN with Residual connections and BatchNorm. Hidden dim $d=32$.
    \item \emph{Decoder:} 3-layer MLP projecting edge embeddings to scalar log-priors.
    \item \emph{Input Featues:} Node coordinates (2D), $k$-NN distances (scalar), pheromone (scalar), heuristic (scalar).
\end{itemize}

\subsection{Complexity Analysis}
We distinguish between neural inference cost and per-iteration ACO costs:

\paragraph{Neural inference (once per outer step).}
The GNN encoder has complexity $O(L \cdot n \cdot K \cdot d)$ for $L=12$ layers, $n$ nodes, $K=32$ candidates per node, and hidden dimension $d=32$.
This runs $H$ times per instance (once per refresh), not every inner iteration.
The edge MLP adds $O(n \cdot K \cdot d)$ per outer step.
Total neural cost per instance: $O(H \cdot L \cdot n \cdot K \cdot d)$.

\paragraph{ACO sampling (each inner iteration).}
Each ant samples at most $M$ edge modifications via sequential next-node selection.
The computational cost per ant is $O(M \cdot K)$ for candidate lookups and probability computation.
With $m$ ants and $T = H \cdot S$ total iterations, sampling costs $O(T \cdot m \cdot M \cdot K)$.

\paragraph{Localized local search.}
Each LS pass scans $O(M)$ checklist nodes, each with $O(K)$ candidate neighbors, yielding $O(m \cdot M \cdot K)$ per inner iteration.

\paragraph{Pheromone update.}
Updating pheromones from the best route is $O(n)$ per inner iteration (sparse update on candidate graph edges).

\paragraph{Summary.}
With default settings ($H=10$, $S=100$, $L=12$, $d=32$, $m=100$, $M=12$, $K=32$), at $n=10$K:
\begin{itemize}
    \item Neural: $\approx 10 \times 12 \times 10\text{K} \times 32 \times 32 \approx 1.2 \times 10^9$ operations
    \item Sampling: $\approx 1000 \times 100 \times 12 \times 32 \approx 3.8 \times 10^7$ operations per call (but simpler ops)
\end{itemize}
In practice, the GNN runs on GPU while sampling runs on CPU (parallelized via OpenMP), so the neural overhead is 5--15\% of total wall-clock time.

\section{CVRP Split Algorithm Details}
\label{app:split}
\emph{Giant Tour Representation:} 
We map the CVRP to a permutation problem by treating the solution as a single sequence of customers. The optimal cost for a given permutation is found by inserting the depot into the sequence to form valid routes.

\emph{Split Algorithm (Prins 2004):}
Let $P = (c_1, c_2, \dots, c_n)$ be the permutation of customers. We construct a directed acyclic graph $H$ where an edge $(i, j)$ exists if the subsequence $(c_{i+1}, \dots, c_j)$ corresponds to a valid route (load $\le$ Capacity). The weight of edge $(i, j)$ is the cost of that route plus the depot connections. Finding the shortest path from node 0 to node $n$ in $H$ yields the optimal route partition.
Let $V_i$ be the cost of the optimal partition for the first $i$ customers. The Bellman equation is:
\[ V_i = \min_{0 \le j < i} \{ V_j + \text{Cost}(c_{j+1}, \dots, c_i) \} \]
where $\text{Cost}(\cdot)$ includes the depot distances.

\emph{Efficient Local Search:}
To maximize efficiency, our Local Search (2-opt) optimizes the \emph{Euclidean length} of the giant tour as a proxy for the CVRP objective. This allows us to evaluate neighbor moves in $O(1)$ using standard distance deltas. The true CVRP cost is computed only once per ant after the LS phase using an efficient \emph{Linear Split} algorithm ($O(N)$) implemented with a monotonic deque \cite{prins2004}, which avoids the quadratic complexity of the naïve Split.

\section{Extended Experimental Plan}
\label{app:experiments}
This appendix provides detailed experimental plans organized by tuning stage and additional robustness studies.

% ======================================================================
% Tuning Pipeline and Hyperparameter Sweeps
% ======================================================================
\subsection{Hyperparameter Sweeps}
\label{app:tuning}
This section documents the hyperparameter sweeps deferred from the main ablation studies.

\paragraph{Solver Hyperparameter Grid (Table~\ref{tab:app-solver}).}
We perform a grid search over core ACO parameters to identify configurations where neural guidance provides the largest improvement over vanilla ACO.
Parameters: $\rho \in \{0.02, 0.05, 0.1, 0.2\}$, $m \in \{50, 100\}$, $K \in \{16, 32\}$, $M \in \{8, 12, 16\}$.

\begin{table}[h]
\centering
\caption{Solver hyperparameter sensitivity. Showing top-5 configurations ranked by neural improvement over vanilla.}
\label{tab:app-solver}
\resizebox{\columnwidth}{!}{%
\begin{tabular}{ccccccc}
\toprule
$\rho$ & $m$ & $K$ & $M$ & Vanilla Gap & Neural Gap & Improv.\ (\%) \\
\midrule
-- & -- & -- & -- & -- & -- & -- \\
-- & -- & -- & -- & -- & -- & -- \\
-- & -- & -- & -- & -- & -- & -- \\
-- & -- & -- & -- & -- & -- & -- \\
-- & -- & -- & -- & -- & -- & -- \\
\bottomrule
\end{tabular}%
}
\end{table}

\paragraph{RL Algorithm Comparison (Table~\ref{tab:app-rl}).}
We compare PPO against REINFORCE across learning rates and PPO-specific hyperparameters.
PPO sweep: $\text{lr} \in \{10^{-6}, 3\times10^{-6}, 10^{-5}, 3\times10^{-5}, 10^{-4}\}$, $\epsilon_{\text{clip}} \in \{0.1, 0.2, 0.3\}$, PPO epochs $\in \{1, 2, 4, 8\}$, with/without advantage normalization.
REINFORCE sweep: $\text{lr} \in \{10^{-6}, 3\times10^{-6}, 10^{-5}, 3\times10^{-5}, 10^{-4}, 3\times10^{-4}\}$.

\begin{table}[h]
\centering
\caption{RL algorithm comparison. Best configuration per algorithm.}
\label{tab:app-rl}
\resizebox{\columnwidth}{!}{%
\begin{tabular}{lccccc}
\toprule
Algorithm & LR & $\epsilon_{\text{clip}}$ & PPO Epochs & Adv Norm & Gap (\%) \\
\midrule
PPO (best) & -- & -- & -- & -- & -- \\
REINFORCE (best) & -- & -- & -- & -- & -- \\
\bottomrule
\end{tabular}%
}
\end{table}

\paragraph{Tuning Pipeline Overview.}
Our experiments follow a 4-stage tuning pipeline implemented in \texttt{tune.py}:

\begin{enumerate}
    \item \textbf{Stage 1 (Solver Tuning):} Grid search over ACO hyperparameters comparing vanilla vs.\ neural-guided performance.
    \item \textbf{Stage 2 (RL Stability):} Sweep over PPO and REINFORCE hyperparameters to find stable training configurations.
    \item \textbf{Stage 3 (Budget Factorization):} Optimize the partition of total iterations into outer/inner iterations.
    \item \textbf{Stage 4 (Component Ablations):} Isolate contributions of individual design choices.
\end{enumerate}

% ======================================================================
% D. Runtime and Scalability (Extended)
% ======================================================================
\subsection{D. Runtime and Scalability (Extended)}

\paragraph{D1. Profiling Breakdown (Table~\ref{tab:profiling-full}).}
Full profiling across all scales with breakdown into: (i) GNN forward pass, (ii) ant sampling, (iii) local search, (iv) pheromone/incumbent updates.

\begin{table}[h]
\centering
\caption{Extended profiling: wall-time breakdown (ms) per outer iteration.}
\label{tab:profiling-full}
\resizebox{\columnwidth}{!}{%
\begin{tabular}{lcccccc}
\toprule
$N$ & GNN & Sampling & LS & Update & Total & Overhead (\%) \\
\midrule
1K & -- & -- & -- & -- & -- & -- \\
5K & -- & -- & -- & -- & -- & -- \\
10K & -- & -- & -- & -- & -- & -- \\
50K & -- & -- & -- & -- & -- & -- \\
100K & -- & -- & -- & -- & -- & -- \\
\bottomrule
\end{tabular}%
}
\end{table}

\paragraph{D2. Scaling Curve Data.}
Raw data for Figure~\ref{fig:scaling}, including:
\begin{itemize}
    \item Gap (\%) at each scale $N \in \{1\text{K}, 5\text{K}, 10\text{K}, 50\text{K}, 100\text{K}\}$
    \item Time budgets used per scale
    \item Comparison against FocusedACO at large scales
\end{itemize}

\paragraph{D3. Time-Quality Tradeoff Curves.}
For each scale, we plot gap vs.\ time budget to show how LGA converges faster than baselines.

% ======================================================================
% E. Robustness and Generalization
% ======================================================================
\subsection{E. Robustness and Generalization}

\paragraph{E1. Distribution Shift (Table~\ref{tab:distribution-shift}).}
We train on uniform random instances and test on out-of-distribution instances:
\begin{itemize}
    \item \emph{Clustered:} 5--10 Gaussian clusters with $\sigma=0.05$
    \item \emph{Mixed:} 50\% uniform, 50\% clustered
    \item \emph{TSPLIB/CVRPLIB:} Real-world benchmark instances
\end{itemize}

\begin{table}[h]
\centering
\caption{E1: Generalization to out-of-distribution instances (trained on uniform).}
\label{tab:distribution-shift}
\resizebox{\columnwidth}{!}{%
\begin{tabular}{lcccc}
\toprule
Distribution & Vanilla Gap & Open-Loop Gap & LGA Gap & Improv.\ (\%) \\
\midrule
Uniform (in-dist) & -- & -- & -- & -- \\
Clustered & -- & -- & -- & -- \\
Mixed & -- & -- & -- & -- \\
TSPLIB & -- & -- & -- & -- \\
\bottomrule
\end{tabular}%
}
\end{table}

\paragraph{E2. Benchmark Evaluation (Table~\ref{tab:benchmarks}).}
Evaluation on standard benchmark sets:
\begin{itemize}
    \item \emph{TSP:} TSPLIB instances (selected large instances)
    \item \emph{CVRP:} Uchoa et al.\ (2017) instances from CVRPLIB
\end{itemize}

\begin{table}[h]
\centering
\caption{E2: Performance on standard benchmarks.}
\label{tab:benchmarks}
\resizebox{\columnwidth}{!}{%
\begin{tabular}{lccccc}
\toprule
Instance & $N$ & Best Known & LGA & Gap (\%) & Time \\
\midrule
\multicolumn{6}{l}{\emph{TSPLIB}} \\
-- & -- & -- & -- & -- & -- \\
\midrule
\multicolumn{6}{l}{\emph{CVRPLIB (Uchoa)}} \\
-- & -- & -- & -- & -- & -- \\
\bottomrule
\end{tabular}%
}
\end{table}

\paragraph{E3. Hyperparameter Robustness (Figure~\ref{fig:robustness}).}
Sensitivity analysis varying one parameter while fixing others:
\begin{itemize}
    \item Candidate list $K \in \{8, 16, 32, 64\}$
    \item Number of ants $m \in \{25, 50, 100, 200\}$
    \item Evaporation rate $\rho \in \{0.02, 0.05, 0.1, 0.2, 0.5\}$
    \item Perturbation limit $M \in \{4, 8, 12, 16, 24\}$
\end{itemize}

% ======================================================================
% F. Training Studies
% ======================================================================
\subsection{F. Training Studies}

\paragraph{F1. Learning Curves (Figure~\ref{fig:learning-curves}).}
Gap vs.\ training steps with $\pm 1$ std bands over 5 seeds.
Compare PPO vs.\ REINFORCE convergence speed and final performance.

\paragraph{F2. Training Stability with Smoothed-MMAS (Figure~\ref{fig:stability}).}
Side-by-side learning curves with and without Smoothed-MMAS.
Expected outcome: standard MMAS shows higher variance and occasional instability.

\paragraph{F3. Size Generalization (Table~\ref{tab:size-gen}).}
Train on small instances, test on larger ones:
\begin{itemize}
    \item Train on $N=1\text{K}$, test on $N \in \{5\text{K}, 10\text{K}\}$
    \item Train on $N=5\text{K}$, test on $N \in \{10\text{K}, 50\text{K}\}$
\end{itemize}

\begin{table}[h]
\centering
\caption{F3: Zero-shot size generalization.}
\label{tab:size-gen}
\small
\begin{tabular}{cccc}
\toprule
Train $N$ & Test $N$ & Gap (\%) & vs.\ Trained-at-Scale \\
\midrule
1K & 5K & -- & -- \\
1K & 10K & -- & -- \\
5K & 10K & -- & -- \\
5K & 50K & -- & -- \\
\bottomrule
\end{tabular}
\end{table}

% ======================================================================
% G. Additional Visualizations
% ======================================================================
\subsection{G. Additional Visualizations}

\paragraph{G1. Solution Visualization.}
Side-by-side tour visualizations for vanilla MFACO vs.\ LGA on representative instances, highlighting structural differences.

\paragraph{G2. Pheromone Heatmaps.}
Evolution of pheromone distribution $\tau$ across outer iterations, showing how LGA shapes the search landscape differently from vanilla ACO.

\paragraph{G3. Prior Logit Visualization.}
Visualization of neural guidance $z_\theta$ at different search stages, showing how the navigator's focus shifts during optimization.

\end{document}
